\documentclass[11pt,a4paper]{article}

\usepackage[utf8]{inputenc}
\usepackage[T1]{fontenc}
\usepackage[french]{babel}
\usepackage{geometry}
\usepackage{xcolor}
\usepackage{graphicx}
\usepackage{float}
\usepackage{array}
\usepackage{longtable}
\usepackage{tabularx}
\usepackage{booktabs}
\usepackage{enumitem}
\usepackage{fancyvrb}
\usepackage{hyperref}
\usepackage{titlesec}

\geometry{margin=2.2cm}
\setlist[itemize]{leftmargin=1.4em}
\setlist[enumerate]{leftmargin=1.6em}
\renewcommand{\arraystretch}{1.18}
\graphicspath{{docs/diagrams/}{diagrams/}}

\definecolor{primary}{HTML}{1D4ED8}
\definecolor{secondary}{HTML}{0F766E}
\definecolor{muted}{HTML}{475569}
\definecolor{soft}{HTML}{F8FAFC}
\definecolor{line}{HTML}{CBD5E1}

\hypersetup{
  colorlinks=true,
  linkcolor=primary,
  urlcolor=primary,
  citecolor=primary,
  pdftitle={Rapport complet - Projet GLPI multi-site},
  pdfauthor={Gournay Lucas, Chapin Florian, De Morais Theo}
}

\titleformat{\section}
  {\Large\bfseries\color{primary}}
  {\thesection}{0.6em}{}

\titleformat{\subsection}
  {\large\bfseries\color{secondary}}
  {\thesubsection}{0.6em}{}

\titleformat{\subsubsection}
  {\normalsize\bfseries\color{muted}}
  {\thesubsubsection}{0.6em}{}

\newcommand{\code}[1]{\texttt{\detokenize{#1}}}
\newcommand{\projecttitle}{Projet GLPI multi-site}
\newcommand{\diagramplaceholder}[2]{%
  \fbox{%
    \begin{minipage}[c][7cm][c]{0.92\textwidth}
      \centering
      \textbf{#2}\\[0.4cm]
      Le schema existe dans le depot sous \code{docs/diagrams/#1.svg}.\\
      Pour l'afficher dans Overleaf, convertir ce fichier en \code{#1.pdf}
      ou \code{#1.png}, puis l'uploader dans \code{docs/diagrams/}.
    \end{minipage}%
  }%
}
\newcommand{\insertdiagram}[2]{%
  \IfFileExists{docs/diagrams/#1.pdf}{%
    \includegraphics[width=\textwidth]{docs/diagrams/#1.pdf}%
  }{%
    \IfFileExists{docs/diagrams/#1.png}{%
      \includegraphics[width=\textwidth]{docs/diagrams/#1.png}%
    }{%
      \IfFileExists{diagrams/#1.pdf}{%
        \includegraphics[width=\textwidth]{diagrams/#1.pdf}%
      }{%
        \IfFileExists{diagrams/#1.png}{%
          \includegraphics[width=\textwidth]{diagrams/#1.png}%
        }{%
          \diagramplaceholder{#1}{#2}%
        }%
      }%
    }%
  }%
}

\begin{document}

\begin{titlepage}
  \centering
  \vspace*{1.2cm}
  {\Huge\bfseries \projecttitle \par}
  \vspace{0.35cm}
  {\Large Rapport complet et detaille\par}
  \vspace{0.3cm}
  {\large Traitement et administration des donnees\par}
  \vfill
  \begin{tabular}{rl}
    Etablissement & CY Tech \\
    Projet & Refonte simplifiee d'une base GLPI sous Oracle XE \\
    Perimetre & Campus de Cergy et Pau \\
    Auteurs & Gournay Lucas, Chapin Florian, De Morais Theo \\
  \end{tabular}
  \vfill
  {\large \today\par}
\end{titlepage}

\tableofcontents
\newpage

\section{Introduction}

\subsection{Contexte general}

GLPI est une solution de gestion de parc informatique et de support. Son
schema relationnel reel est tres riche, car il couvre un grand nombre de
modules : inventaire, tickets, contrats, reseau, utilisateurs, droits,
plugins et historique. Cette richesse devient cependant difficile a exploiter
dans le cadre d'un mini-projet de base de donnees repartie : le modele initial
comporte beaucoup de tables, peu de relations explicites faciles a presenter et
plusieurs structures tres proches pour des objets materiels differents.

Le projet propose donc une refonte simplifiee d'une base inspiree de GLPI pour
deux campus, Cergy et Pau. L'objectif n'est pas de reproduire toute la base
GLPI, mais de construire un noyau coherent, contraint, explicable et
demonstrable sous Oracle XE. La conception conserve les notions importantes :
organisation multi-sites, utilisateurs, profils, groupes, inventaire, tickets,
reseau, audit, archivage, optimisation physique et simulation BDDR.

\subsection{Objectifs du projet}

Le projet poursuit cinq objectifs principaux, repris du plan de presentation :

\begin{enumerate}
  \item realiser un reverse engineering du modele GLPI pour identifier les
  entites utiles au contexte pedagogique ;
  \item produire un modele simplifie, lisible et justifie par un MCD et un MLD ;
  \item proposer une architecture de base de donnees repartie entre Cergy et Pau ;
  \item ajouter une logique metier PL/SQL garantissant l'integrite et les
  traitements principaux ;
  \item optimiser les acces avec tablespaces, clusters, index et mesures de
  performance.
\end{enumerate}

\subsection{Plan suivi}

Le rapport suit le plan de la presentation :

\begin{longtable}{p{4cm}p{10cm}}
\toprule
\textbf{Partie de presentation} & \textbf{Traitement dans le rapport} \\
\midrule
Contexte et objectifs & Perimetre, limites de GLPI et objectifs Oracle. \\
Reverse engineering et modelisation & Choix de simplification, MCD, MLD et contraintes. \\
Architecture repartie & Fragmentation horizontale, schemas Cergy/Pau, DB links et vues globales. \\
Optimisation physique & Tablespaces, clusters, index B-tree, composites, fonctionnels et bitmap. \\
Logique metier PL/SQL & Procedures, fonctions, triggers, audit et securite. \\
Analyse des performances & Protocole de benchmark et interpretation attendue. \\
\bottomrule
\end{longtable}

\section{Analyse du probleme initial}

\subsection{Complexite du modele GLPI}

Le schema GLPI original contient plusieurs centaines de tables. Beaucoup de
modules ne sont pas necessaires pour le perimetre du projet : contrats,
licences, notifications, calendriers avances, plugins, historique detaille,
relations reseau tres fines, etc. Conserver tout ce schema aurait rendu la base
difficile a expliquer et a distribuer proprement.

Le premier enjeu a donc ete de separer le coeur fonctionnel du bruit technique.
Le coeur retenu correspond aux besoins suivants :

\begin{itemize}
  \item connaitre les sites et localisations physiques ;
  \item gerer les utilisateurs et leurs droits principaux ;
  \item inventorier les equipements informatiques ;
  \item suivre les incidents ou demandes lies a un materiel ;
  \item associer les materiels a des ports reseau et a des adresses IP ;
  \item tracer les modifications importantes.
\end{itemize}

\subsection{Limites d'un modele trop proche de GLPI}

Un modele trop proche de GLPI aurait pose plusieurs difficultes :

\begin{itemize}
  \item multiplication des tables de materiel, par exemple ordinateurs,
  ecrans, imprimantes, telephones, peripheriques et equipements reseau ;
  \item nombreuses associations polymorphes, difficiles a representer avec des
  cles etrangeres simples ;
  \item vues et requetes plus lourdes, car l'inventaire complet doit etre
  reconstruit par unions ou jointures multiples ;
  \item distribution moins lisible, car chaque famille d'objets aurait son
  propre fragment ;
  \item demonstration plus fragile, puisque le jury doit comprendre rapidement
  le modele et ses choix.
\end{itemize}

\subsection{Strategie de simplification}

La simplification principale consiste a centraliser tous les equipements dans
une table unique \code{assets}. La colonne \code{asset_type} indique la nature du
materiel :

\begin{center}
\begin{tabular}{ll}
\toprule
\textbf{Type} & \textbf{Sens fonctionnel} \\
\midrule
\code{COMPUTER} & Ordinateur fixe ou portable. \\
\code{MONITOR} & Ecran. \\
\code{PERIPHERAL} & Peripherique general. \\
\code{PRINTER} & Imprimante. \\
\code{PHONE} & Telephone. \\
\code{NETWORK_EQUIPMENT} & Switch, borne Wi-Fi, routeur ou equipement reseau. \\
\bottomrule
\end{tabular}
\end{center}

Ce choix remplace plusieurs tables quasiment identiques par une structure
commune. Il reduit les jointures, simplifie les vues metier et rend l'ajout
d'un nouveau type de materiel possible sans modifier le schema.

\section{Reverse engineering et modelisation}

\subsection{Domaines fonctionnels conserves}

Le modele final contient 19 tables, dont 17 tables fonctionnelles et 2 tables
systeme. Elles sont reparties en six domaines :

\begin{longtable}{p{3.5cm}p{10.5cm}}
\toprule
\textbf{Domaine} & \textbf{Tables} \\
\midrule
Organisation & \code{sites}, \code{locations}. \\
Referentiels & \code{manufacturers}, \code{states}, \code{ticket_categories}. \\
Utilisateurs et droits & \code{users}, \code{profiles}, \code{profiles_users}, \code{groups}, \code{groups_users}. \\
Inventaire & \code{assets}. \\
Support & \code{tickets}, \code{ticket_users}, \code{ticket_followups}. \\
Reseau & \code{network_ports}, \code{ip_networks}, \code{ip_addresses}. \\
Systeme & \code{audit_log}, \code{archives_materiel}. \\
\bottomrule
\end{longtable}

\subsection{Comparaison avant et apres simplification}

\begin{longtable}{p{6.5cm}p{7.5cm}}
\toprule
\textbf{Probleme du modele source} & \textbf{Choix du modele final} \\
\midrule
Tables distinctes pour chaque type de materiel. & Une seule table \code{assets}, typee par \code{asset_type}. \\
Relations polymorphes dans les tickets et les ports reseau. & Une seule cle etrangere \code{asset_id}. \\
Nombreux referentiels utilisateurs secondaires. & Conservation des profils, groupes et affectations essentielles. \\
Modele reseau tres fin. & Conservation des ports, sous-reseaux, VLAN et adresses IP. \\
Historique disperse. & Tables systeme dediees pour l'audit et l'archivage. \\
\bottomrule
\end{longtable}

\subsection{MCD du projet}

Le MCD represente les entites retenues et leurs associations principales. Il
met volontairement en avant la table centrale \code{assets}, reliee aux sites,
aux utilisateurs, aux localisations, aux fabricants, aux etats, aux tickets et
aux ports reseau.

\begin{figure}[H]
  \centering
  \insertdiagram{mcd}{MCD simplifie de la base GLPI multi-site}
  \caption{MCD simplifie de la base GLPI multi-site.}
  \label{fig:mcd}
\end{figure}

La lecture du MCD se fait autour de quatre axes :

\begin{itemize}
  \item un site possede des localisations, des utilisateurs, des groupes, des
  equipements, des tickets et des elements reseau ;
  \item un utilisateur peut posseder un materiel, etre technicien responsable,
  demander un ticket ou intervenir dans un suivi ;
  \item un materiel peut avoir plusieurs tickets et plusieurs ports reseau ;
  \item un ticket possede des affectations de techniciens et des suivis.
\end{itemize}

\subsection{Justification des cardinalites}

\begin{longtable}{p{6cm}p{3cm}p{5cm}}
\toprule
\textbf{Relation} & \textbf{Cardinalite} & \textbf{Justification} \\
\midrule
\code{sites -> locations} & 1,N & Un campus ou service peut contenir plusieurs batiments ou salles. \\
\code{sites -> users} & 1,N & Un utilisateur appartient a un site principal. \\
\code{sites -> assets} & 1,N & Chaque equipement est rattache a un site proprietaire. \\
\code{assets -> tickets} & 1,N & Un materiel peut generer plusieurs demandes support. \\
\code{assets -> network_ports} & 1,N & Un equipement peut exposer plusieurs interfaces reseau. \\
\code{network_ports -> ip_addresses} & 1,N & Un port peut recevoir une ou plusieurs adresses selon le contexte. \\
\code{tickets -> ticket_followups} & 1,N & Un ticket contient un historique de commentaires. \\
\code{users -> profiles_users} & 1,N & Un utilisateur peut avoir plusieurs profils selon le site. \\
\bottomrule
\end{longtable}

\section{Modele logique de donnees}

\subsection{Principes du MLD}

Le MLD traduit le MCD en tables relationnelles Oracle. Les choix structurants
sont les suivants :

\begin{itemize}
  \item chaque table possede une cle primaire numerique \code{id} ;
  \item les relations sont exprimees par des cles etrangeres lisibles ;
  \item les tables de liaison portent les associations N,N ;
  \item les doublons fonctionnels sont evites par des contraintes
  \code{UNIQUE} ;
  \item les suppressions dependantes sont encadrees avec \code{ON DELETE CASCADE}
  ou \code{ON DELETE SET NULL} selon le sens metier.
\end{itemize}

\subsection{Organisation}

\begin{Verbatim}[fontsize=\small, frame=single]
sites(
  id PK,
  name,
  parent_site_id FK -> sites.id,
  full_name,
  site_code,
  created_at,
  updated_at
)

locations(
  id PK,
  site_id FK -> sites.id,
  name,
  parent_location_id FK -> locations.id,
  full_name,
  building,
  room,
  created_at,
  updated_at
)
\end{Verbatim}

\code{sites} porte la hierarchie logique et le code de distribution
\code{site_code}. \code{locations} represente les lieux physiques, comme un
batiment ou une salle.

\subsection{Referentiels}

\begin{Verbatim}[fontsize=\small, frame=single]
manufacturers(id PK, name UNIQUE)
states(id PK, name UNIQUE)
ticket_categories(id PK, name UNIQUE, description)
\end{Verbatim}

Ces tables sont stables et peu volumineuses. Elles sont donc de bonnes
candidates a la replication logique dans les deux schemas de simulation BDDR.

\subsection{Utilisateurs, profils et groupes}

\begin{Verbatim}[fontsize=\small, frame=single]
users(
  id PK,
  login UNIQUE,
  last_name,
  first_name,
  email,
  phone,
  site_id FK -> sites.id,
  location_id FK -> locations.id,
  supervisor_user_id FK -> users.id,
  is_active,
  created_at,
  updated_at
)

profiles(id PK, name UNIQUE, interface, is_default)

profiles_users(
  id PK,
  user_id FK -> users.id,
  profile_id FK -> profiles.id,
  site_id FK -> sites.id,
  is_recursive,
  UNIQUE(user_id, profile_id, site_id)
)

groups(
  id PK,
  site_id FK -> sites.id,
  name,
  parent_group_id FK -> groups.id,
  full_name,
  created_at,
  updated_at
)

groups_users(
  id PK,
  user_id FK -> users.id,
  group_id FK -> groups.id,
  is_manager,
  UNIQUE(user_id, group_id)
)
\end{Verbatim}

Le couple \code{profiles} et \code{profiles_users} reprend l'idee GLPI
d'affectation de profils selon un contexte. Les groupes servent principalement
a representer les equipes support locales.

\subsection{Inventaire}

\begin{Verbatim}[fontsize=\small, frame=single]
assets(
  id PK,
  site_id FK -> sites.id,
  asset_type,
  name,
  serial_number,
  owner_user_id FK -> users.id,
  technician_user_id FK -> users.id,
  location_id FK -> locations.id,
  manufacturer_id FK -> manufacturers.id,
  state_id FK -> states.id,
  created_at,
  updated_at,
  UNIQUE(site_id, serial_number)
)
\end{Verbatim}

\code{assets} est la table centrale du projet. La contrainte
\code{UNIQUE(site_id, serial_number)} evite deux numeros de serie identiques
sur un meme site, tout en laissant la possibilite theorique qu'un meme serial
soit reutilise dans un autre fragment si le cas metier le permet.

\subsection{Support}

\begin{Verbatim}[fontsize=\small, frame=single]
tickets(
  id PK,
  site_id FK -> sites.id,
  asset_id FK -> assets.id ON DELETE CASCADE,
  title,
  description,
  status,
  priority,
  requester_user_id FK -> users.id,
  assigned_group_id FK -> groups.id,
  category_id FK -> ticket_categories.id,
  resolution,
  created_at,
  updated_at,
  assigned_at,
  resolved_at,
  closed_at
)

ticket_users(
  id PK,
  ticket_id FK -> tickets.id ON DELETE CASCADE,
  user_id FK -> users.id,
  assigned_by_user_id FK -> users.id,
  assigned_at,
  UNIQUE(ticket_id, user_id)
)

ticket_followups(
  id PK,
  ticket_id FK -> tickets.id ON DELETE CASCADE,
  user_id FK -> users.id,
  content,
  created_at
)
\end{Verbatim}

Le support est volontairement centre sur un materiel. Un ticket est associe a
un \code{asset_id}, ce qui simplifie l'analyse des incidents et la creation des
vues metier.

\subsection{Reseau}

\begin{Verbatim}[fontsize=\small, frame=single]
network_ports(
  id PK,
  site_id FK -> sites.id,
  asset_id FK -> assets.id ON DELETE CASCADE,
  port_name,
  mac_address,
  port_type,
  created_at,
  updated_at
)

ip_networks(
  id PK,
  site_id FK -> sites.id,
  network_name,
  network_address,
  subnet_mask,
  gateway_address,
  vlan_name,
  vlan_tag,
  created_at,
  updated_at,
  UNIQUE(site_id, network_address, subnet_mask)
)

ip_addresses(
  id PK,
  site_id FK -> sites.id,
  network_port_id FK -> network_ports.id ON DELETE SET NULL,
  ip_network_id FK -> ip_networks.id,
  ip_address,
  created_at,
  updated_at
)
\end{Verbatim}

Le reseau est limite aux elements utiles pour une demonstration : ports, sous-
reseaux, VLAN et adresses IP. Cette granularite suffit pour montrer des
jointures et des statistiques sans introduire un modele reseau trop lourd.

\subsection{Tables systeme}

\begin{Verbatim}[fontsize=\small, frame=single]
audit_log(
  id PK,
  table_name,
  row_id,
  action,
  old_data,
  new_data,
  changed_by,
  changed_at
)

archives_materiel(
  id PK,
  original_table,
  original_id,
  archived_data,
  archived_by,
  archived_at
)
\end{Verbatim}

\code{audit_log} conserve les modifications significatives. \code{archives_materiel}
garde une trace JSON d'un materiel avant suppression definitive.

\section{Architecture repartie BDDR}

\subsection{Principe general}

La BDDR est simulee localement sous Oracle XE avec deux schemas :
\code{GLPI_CERGY} et \code{GLPI_PAU}. Chaque schema expose le fragment logique
de son site. Cette simulation garde les principes d'une base repartie tout en
restant executable sur un seul poste.

\begin{figure}[H]
  \centering
  \insertdiagram{architecture}{Architecture Oracle XE distribuee entre Cergy et Pau}
  \caption{Architecture Oracle XE distribuee entre Cergy et Pau.}
  \label{fig:architecture}
\end{figure}

\subsection{Fragmentation horizontale}

La fragmentation repose sur \code{sites.site_code}. Les donnees operationnelles
sont decoupees selon le site :

\begin{longtable}{p{5cm}p{8.5cm}}
\toprule
\textbf{Famille de donnees} & \textbf{Strategie} \\
\midrule
Sites et localisations & Chaque schema expose les lignes de son site. \\
Utilisateurs et groupes & Fragmentation par site de rattachement. \\
Assets & Fragmentation par \code{assets.site_id}. \\
Tickets et suivis & Fragmentation selon le site du ticket et du materiel. \\
Reseau & Ports, sous-reseaux et IP filtres par site. \\
Referentiels & Exposition commune dans les deux schemas. \\
\bottomrule
\end{longtable}

Ce choix est naturel : Cergy heberge les equipements, utilisateurs, tickets et
adresses de Cergy ; Pau heberge ceux de Pau. Les referentiels sont repliques
car ils changent rarement et doivent rester disponibles des deux cotes.

\subsection{Schemas, vues et synonyms}

Le script \code{07_bddr.sql} cree les utilisateurs \code{glpi_cergy} et
\code{glpi_pau}, puis expose des vues locales dans chaque schema. Les vues
permettent de masquer le schema source et d'offrir une interface logique comme
si chaque site possedait ses propres tables.

Les database links \code{lien_pau} et \code{lien_cergy} permettent ensuite
d'interroger le site distant. Des synonymes, par exemple \code{assets_pau} ou
\code{assets_cergy}, masquent la syntaxe distante \code{@lien_pau} ou
\code{@lien_cergy}.

\subsection{Vues globales}

Les vues globales fournissent une vision consolidee :

\begin{longtable}{p{4.5cm}p{9.5cm}}
\toprule
\textbf{Vue globale} & \textbf{Role} \\
\midrule
\code{V_ASSETS_GLOBAL} & Inventaire Cergy + Pau avec indication du fragment. \\
\code{V_USERS_GLOBAL} & Utilisateurs des deux sites. \\
\code{V_STATS_GLOBAL} & Statistiques agregees par site. \\
\code{V_TICKETS_GLOBAL} & Tickets consolides. \\
\bottomrule
\end{longtable}

Cette approche montre le principe classique d'une BDDR : les applications
locales travaillent sur leur fragment, tandis que les besoins de reporting
passent par des vues globales.

\subsection{Transfert inter-sites}

La procedure \code{SP_TRANSFERT_MATERIEL} represente le traitement cle de la
BDDR. Son role est de deplacer un materiel d'un site vers l'autre. Dans une
situation locale, le changement peut se limiter a la mise a jour de
\code{assets.site_id}. Dans un cas inter-sites, le transfert doit tenir compte
du fragment distant :

\begin{enumerate}
  \item identifier le site courant du materiel ;
  \item identifier le site cible ;
  \item recopier ou mettre a jour le materiel dans le fragment distant ;
  \item deplacer les dependances utiles, notamment ports reseau, adresses IP,
  tickets, affectations et suivis ;
  \item supprimer l'ancien fragment local lorsque le transfert est termine.
\end{enumerate}

Cette logique illustre le probleme central des BDDR : une operation metier
simple en apparence peut impliquer plusieurs objets repartis et doit preserver
l'integrite referentielle.

\section{Integrite, securite et droits}

\subsection{Contraintes relationnelles}

L'integrite est d'abord assuree par le schema relationnel :

\begin{itemize}
  \item les cles primaires identifient chaque ligne ;
  \item les cles etrangeres imposent la coherence entre domaines ;
  \item les contraintes \code{UNIQUE} evitent les doublons fonctionnels ;
  \item les contraintes \code{CHECK} encadrent les valeurs enumerees comme
  \code{asset_type}, \code{status} ou \code{priority}.
\end{itemize}

La contrainte \code{UNIQUE(site_id, serial_number)} sur \code{assets} remplace
avantageusement un trigger de verification de serial, car elle evite les
problemes de table mutante et laisse Oracle garantir l'unicite directement.

\subsection{Roles Oracle}

Le projet definit quatre roles principaux :

\begin{longtable}{p{4cm}p{10cm}}
\toprule
\textbf{Role} & \textbf{Responsabilite} \\
\midrule
\code{ROLE_ADMIN} & Administration complete du schema et des objets. \\
\code{ROLE_TECHNICIEN} & Consultation et actions sur le support, l'inventaire et le reseau. \\
\code{ROLE_CONSULTANT} & Lecture des vues metier. \\
\code{ROLE_MANAGER_SITE} & Suivi des donnees et indicateurs du site. \\
\bottomrule
\end{longtable}

Ces roles separents les usages : administration, intervention technique,
consultation et pilotage local.

\subsection{Audit et archivage}

L'audit est realise par \code{TRG_AUDIT_ASSETS}, qui journalise les operations
sur les materiels. L'archivage est realise par
\code{TRG_ARCHIVE_ASSET_DELETE}, qui conserve l'etat final d'un materiel avant
suppression.

Ce mecanisme repond a deux besoins :

\begin{itemize}
  \item garder une trace des modifications importantes ;
  \item pouvoir expliquer ce qui est arrive a un materiel supprime ou transfere.
\end{itemize}

\section{Optimisation physique}

\subsection{Tablespaces}

Le projet separe les donnees dans plusieurs tablespaces :

\begin{longtable}{p{4cm}p{9.5cm}}
\toprule
\textbf{Tablespace} & \textbf{Utilisation} \\
\midrule
\code{TS_MATERIEL} & Inventaire et donnees associees aux assets. \\
\code{TS_UTILISATEURS} & Utilisateurs, profils et groupes. \\
\code{TS_RESEAU} & Ports, sous-reseaux et adresses IP. \\
\code{TS_SUPPORT} & Tickets et suivis support. \\
\code{TS_INDEX} & Index et structures d'optimisation. \\
\bottomrule
\end{longtable}

Cette separation rend l'organisation physique plus claire et permet de
presenter une logique d'administration Oracle : stockage par domaine,
surveillance des volumes et isolation des index.

\subsection{Clusters}

Les clusters regroupent physiquement des lignes souvent lues ensemble. Le
projet utilise trois clusters :

\begin{longtable}{p{5cm}p{3cm}p{6cm}}
\toprule
\textbf{Cluster} & \textbf{Cle} & \textbf{Tables concernees} \\
\midrule
\code{cluster_user_rights} & \code{user_id} & \code{users}, \code{profiles_users}, \code{groups_users}. \\
\code{cluster_site_structure} & \code{site_id} & \code{sites}, \code{locations}, \code{groups}, \code{ip_networks}. \\
\code{cluster_network_port_ips} & \code{network_port_id} & \code{network_ports}, \code{ip_addresses}. \\
\bottomrule
\end{longtable}

Le cluster \code{cluster_user_rights} est particulierement utile pour expliquer
l'objectif : lorsqu'on consulte un utilisateur, ses profils et groupes sont des
donnees frequemment accedees ensemble. Les regrouper physiquement peut reduire
le nombre de lectures disque.

\subsection{Index B-tree et composites}

Les index B-tree couvrent les cles etrangeres et les recherches frequentes.
Les index composites ciblent les requetes metier :

\begin{itemize}
  \item \code{idx_assets_site_type} pour l'inventaire par site et type ;
  \item \code{idx_assets_site_state} pour l'inventaire par site et etat ;
  \item \code{idx_assets_site_owner} pour retrouver le materiel d'un utilisateur ;
  \item \code{idx_ticket_site_status_priority} pour suivre les tickets ouverts
  ou prioritaires par site ;
  \item \code{idx_ticket_asset_status} pour analyser les tickets d'un materiel ;
  \item \code{idx_network_site_address} pour les sous-reseaux d'un site.
\end{itemize}

\subsection{Index fonctionnels et bitmap}

Les index fonctionnels facilitent les recherches insensibles a la casse :
\code{idx_assets_name_upper}, \code{idx_assets_serial_upper},
\code{idx_users_login_upper} et \code{idx_users_email_upper}.

Les index bitmap sont utilises sur des colonnes a faible cardinalite, comme
\code{asset_type}, \code{is_active}, \code{priority} et \code{port_type}. Ils
sont interessants pour les requetes analytiques, notamment sur de gros volumes,
mais doivent etre presentes avec prudence dans un contexte tres transactionnel.

\section{Vues metier}

\subsection{Role des vues}

Les vues servent a fournir une abstraction metier au-dessus du modele
relationnel. Elles evitent de demander aux utilisateurs ou au jury de retenir
toutes les jointures.

\begin{longtable}{p{4.8cm}p{9.2cm}}
\toprule
\textbf{Vue} & \textbf{Utilite} \\
\midrule
\code{V_INVENTAIRE_COMPLET} & Inventaire enrichi avec site, localisation, proprietaire, technicien, fabricant et etat. \\
\code{V_MATERIEL_PAR_SITE} & Comptage des materiels par site et par type. \\
\code{V_UTILISATEURS_PROFILS} & Utilisateurs avec profils et site d'application. \\
\code{V_TOPOLOGIE_RESEAU} & Ports, adresses IP, sous-reseaux et VLAN. \\
\code{V_STATISTIQUES_SITE} & Indicateurs synthetiques par site. \\
\code{V_MATERIEL_RECENT} & Materiels modifies recemment. \\
\code{V_TICKETS_SUPPORT} & Tickets enrichis avec demandeur, categorie, groupe et materiel. \\
\bottomrule
\end{longtable}

\subsection{Exemple : inventaire complet}

\code{V_INVENTAIRE_COMPLET} est la vue la plus importante pour la
demonstration. Elle montre directement l'interet de la table \code{assets} :
une seule vue suffit pour lister l'ensemble du parc, quel que soit le type de
materiel.

\section{Logique metier PL/SQL}

\subsection{Procedures}

\begin{longtable}{p{5cm}p{9cm}}
\toprule
\textbf{Procedure} & \textbf{Role} \\
\midrule
\code{SP_TRANSFERT_MATERIEL} & Transfert local ou inter-sites d'un materiel. \\
\code{SP_AFFECTER_PROFIL} & Affectation ou mise a jour d'un profil utilisateur. \\
\code{SP_CREER_TICKET_MATERIEL} & Creation d'un ticket sur un asset existant du site. \\
\code{SP_ASSIGNER_TECH_TICKET} & Affectation d'un technicien a un ticket. \\
\code{SP_INVENTAIRE_SITE} & Affichage synthetique de l'inventaire d'un site. \\
\code{SP_NETTOYER_ARCHIVES} & Suppression des archives trop anciennes. \\
\code{SP_GENERER_JEU_DE_TEST} & Generation d'un jeu de donnees representatif. \\
\code{SP_REPLIQUER_REFERENTIELS} & Procedure de demonstration liee a la BDDR. \\
\bottomrule
\end{longtable}

\subsection{Fonctions}

\begin{longtable}{p{5cm}p{9cm}}
\toprule
\textbf{Fonction} & \textbf{Role} \\
\midrule
\code{FN_COMPTER_MATERIEL_SITE} & Compte les assets d'un site. \\
\code{FN_CALCULER_TAUX_UTILISATION} & Calcule le taux d'utilisation reseau d'un asset. \\
\code{FN_OBTENIR_SITE_ENTITE} & Retrouve le code site associe a une entite. \\
\code{FN_VERIFIER_QUOTA_MATERIEL} & Verifie si un quota de materiel est depasse. \\
\bottomrule
\end{longtable}

\subsection{Triggers}

\begin{longtable}{p{5.2cm}p{3.2cm}p{5.5cm}}
\toprule
\textbf{Trigger} & \textbf{Table} & \textbf{Role} \\
\midrule
\code{TRG_AUTO_DATE_MOD_ASSETS} & \code{assets} & Mise a jour de \code{updated_at}. \\
\code{TRG_AUTO_DATE_MOD_USERS} & \code{users} & Mise a jour de \code{updated_at}. \\
\code{TRG_AUTO_DATE_MOD_NETPORTS} & \code{network_ports} & Mise a jour de \code{updated_at}. \\
\code{TRG_AUTO_DATE_MOD_TICKETS} & \code{tickets} & Gestion des dates du workflow ticket. \\
\code{TRG_CHECK_TICKET_USER_TECH} & \code{ticket_users} & Controle qu'un technicien affecte possede le bon profil. \\
\code{TRG_AUDIT_ASSETS} & \code{assets} & Journalisation des modifications d'inventaire. \\
\code{TRG_ARCHIVE_ASSET_DELETE} & \code{assets} & Archivage JSON avant suppression. \\
\code{TRG_CASCADE_SITE_CODE} & \code{sites} & Propagation d'un changement de code site. \\
\bottomrule
\end{longtable}

\subsection{Interet pedagogique du PL/SQL}

Le PL/SQL montre que la base ne se limite pas a stocker des donnees. Elle porte
aussi des regles :

\begin{itemize}
  \item automatiser les dates de modification ;
  \item empecher certaines affectations incoherentes ;
  \item produire des indicateurs ;
  \item encapsuler les operations sensibles comme la creation d'un ticket ou le
  transfert d'un materiel.
\end{itemize}

\section{Jeu de donnees et demonstration}

\subsection{Generation des donnees}

Le script \code{09_test_data.sql} cree la procedure
\code{SP_GENERER_JEU_DE_TEST}. Elle genere un volume de donnees suffisant pour
tester les vues et les performances :

\begin{itemize}
  \item sites Cergy et Pau ;
  \item localisations ;
  \item utilisateurs ;
  \item profils, groupes et affectations ;
  \item assets de plusieurs types ;
  \item tickets, techniciens et suivis ;
  \item ports reseau, sous-reseaux et adresses IP.
\end{itemize}

Le jeu de test permet de montrer les requetes sans saisir manuellement des
dizaines de lignes pendant la soutenance.

\subsection{Scenario de demonstration}

Une demonstration efficace peut suivre cet ordre :

\begin{enumerate}
  \item presenter le MCD et expliquer la centralite de \code{assets} ;
  \item montrer le MLD et les contraintes principales ;
  \item executer les vues metier, notamment \code{V_INVENTAIRE_COMPLET} et
  \code{V_TICKETS_SUPPORT} ;
  \item creer un ticket avec \code{SP_CREER_TICKET_MATERIEL} ;
  \item afficher l'audit ou les dates modifiees par trigger ;
  \item montrer les schemas \code{GLPI_CERGY} et \code{GLPI_PAU} ;
  \item interroger une vue globale BDDR ;
  \item expliquer le benchmark avec index visibles puis invisibles.
\end{enumerate}

\section{Analyse des performances}

\subsection{Protocole}

Le protocole presente dans les slides repose sur l'idee suivante :

\begin{enumerate}
  \item injecter un volume important de donnees avec \code{SP_GENERER_JEU_DE_TEST} ;
  \item vider ou neutraliser autant que possible les effets du cache Oracle
  pour obtenir des mesures comparables ;
  \item executer les requetes representatives ;
  \item comparer les resultats avec index visibles et index invisibles ;
  \item stocker les mesures dans \code{benchmark_results}.
\end{enumerate}

\subsection{Requetes representatives}

Les benchmarks visent les usages principaux :

\begin{itemize}
  \item recherche d'un asset par nom ;
  \item recherche par numero de serie ;
  \item inventaire d'un site par type ;
  \item tickets ouverts par site, statut et priorite ;
  \item topologie reseau d'un materiel ;
  \item droits d'un utilisateur via profils et groupes.
\end{itemize}

\subsection{Resultats de test}

Le tableau suivant reprend les mesures obtenues pendant les tests. Chaque
requete a ete executee dans deux scenarios : avec les index actifs, puis sans
les index. Les temps sont exprimes en millisecondes.

\begin{longtable}{p{1.3cm}p{6.4cm}p{2.6cm}p{2.6cm}}
\toprule
\textbf{ID} & \textbf{Requete testee} & \textbf{Avec index} & \textbf{Sans index} \\
\midrule
\code{Q1} & Recherche materiel par nom & 10 ms & 30 ms \\
\code{Q2} & Inventaire par site et categorie & 50 ms & 260 ms \\
\code{Q3} & Tickets ouverts par priorite & 20 ms & 30 ms \\
\code{Q4} & Recherche materiel par numero de serie & 0 ms & 40 ms \\
\code{Q5} & Reseau des appareils & 190 ms & 220 ms \\
\code{Q6} & Utilisateurs avec profils & 10 ms & 30 ms \\
\bottomrule
\end{longtable}

Ces resultats confirment que les index sont surtout efficaces sur les recherches
selectives et les inventaires filtres. Le gain le plus visible concerne
\code{Q2}, ou l'inventaire par site et categorie passe de 260 ms sans index a
50 ms avec index. Les requetes reseau, comme \code{Q5}, profitent moins des
index car elles reposent davantage sur des jointures et sur un volume de lignes
plus important.

\subsection{Interpretation attendue}

Les index doivent ameliorer surtout les recherches selectives et les filtres
frequents. Les index composites sont utiles quand les colonnes sont utilisees
ensemble, par exemple \code{site_id}, \code{status} et \code{priority} pour les
tickets. Les clusters deviennent interessants lorsque plusieurs tables sont
lues ensemble de maniere repetee, comme les droits utilisateurs.

Le resultat attendu n'est pas seulement un temps plus faible. Il faut aussi
montrer que chaque index correspond a un usage metier. Un index inutile coute de
l'espace et ralentit les insertions ; un bon index accelere une requete
importante et justifie son cout de maintenance.

\section{Limites et ameliorations possibles}

\subsection{Limites du projet}

Le projet reste une simplification. Certaines fonctionnalites completes de GLPI
ne sont pas modelisees :

\begin{itemize}
  \item contrats, fournisseurs, licences logicielles et budgets ;
  \item historique complet de tous les objets ;
  \item droits fins comparables au moteur GLPI reel ;
  \item synchronisation BDDR reelle entre deux serveurs physiques ;
  \item resolution automatique des conflits inter-sites.
\end{itemize}

\subsection{Ameliorations}

Plusieurs evolutions sont possibles :

\begin{itemize}
  \item separer la demonstration locale et une version multi-instances Oracle ;
  \item ajouter des materialized views pour certains rapports globaux ;
  \item enrichir l'audit avec le contexte applicatif ;
  \item ajouter une table de mouvements d'inventaire pour tracer les transferts ;
  \item comparer les plans d'execution avant et apres chaque index ;
  \item automatiser la generation d'un rapport de benchmark.
\end{itemize}

\section{Conclusion}

Le projet aboutit a une base Oracle coherente et defendable. La refonte
simplifiee conserve les idees essentielles de GLPI tout en supprimant la
complexite non necessaire au cadre pedagogique. La table \code{assets} rend
l'inventaire lisible, les vues fournissent une abstraction metier, la BDDR
illustre la fragmentation Cergy/Pau, et le PL/SQL ajoute les regles de gestion
attendues.

La proposition est donc pertinente pour une soutenance : elle montre a la fois
la modelisation, l'administration Oracle, la distribution des donnees, la
securite, l'integrite, les traitements stockes et l'optimisation des
performances. Le modele n'est pas une copie exhaustive de GLPI ; c'est une
version volontairement recentree, plus facile a expliquer et a tester.

\appendix

\section{Annexe A - Dictionnaire synthetique des tables}

\begin{longtable}{p{3.6cm}p{10.4cm}}
\toprule
\textbf{Table} & \textbf{Description} \\
\midrule
\code{sites} & Sites, campus et structures logiques. \\
\code{locations} & Emplacements physiques rattaches aux sites. \\
\code{manufacturers} & Fabricants des equipements. \\
\code{states} & Etats fonctionnels des equipements. \\
\code{users} & Utilisateurs, demandeurs, techniciens et responsables. \\
\code{profiles} & Profils applicatifs. \\
\code{profiles_users} & Affectation des profils aux utilisateurs selon le site. \\
\code{groups} & Groupes de support ou d'organisation. \\
\code{groups_users} & Membres des groupes. \\
\code{assets} & Inventaire centralise des materiels. \\
\code{ticket_categories} & Categories de tickets. \\
\code{tickets} & Demandes et incidents. \\
\code{ticket_users} & Techniciens affectes aux tickets. \\
\code{ticket_followups} & Historique des suivis de tickets. \\
\code{network_ports} & Ports reseau des assets. \\
\code{ip_networks} & Sous-reseaux IP et VLAN. \\
\code{ip_addresses} & Adresses IP rattachees aux ports. \\
\code{audit_log} & Journal d'audit. \\
\code{archives_materiel} & Archives JSON des materiels supprimes. \\
\bottomrule
\end{longtable}

\section{Annexe B - Ordre d'execution des scripts}

\begin{Verbatim}[fontsize=\small, frame=single]
01_tablespaces.sql
02_schema_tables.sql
04_clusters_indexes.sql
03_users_roles.sql
05_views.sql
06_plsql/triggers.sql
06_plsql/procedures.sql
06_plsql/functions.sql
06_plsql/cursors.sql
09_test_data.sql
08_query_plans.sql
10_benchmark.sql
07_bddr.sql
\end{Verbatim}

\end{document}
